\chapter{La circunferencia}

\begin{center}\begin{minipage}{.88\textwidth}\small\sffamily
\textbf{Propósito.} Modelar circunferencias a partir de su centro, radio,
ecuación y relaciones con puntos y rectas.
\end{minipage}\end{center}

\begin{definition}{Circunferencia}{circunferencia}
Es el lugar geométrico de los puntos $P(x,y)$ cuya distancia a un punto fijo
$C(h,k)$, llamado centro, es una constante positiva $r$, llamada radio.
\end{definition}

De $d(P,C)=r$ se obtiene
\[
\sqrt{(x-h)^2+(y-k)^2}=r.
\]
Al elevar al cuadrado aparece su ecuación ordinaria.

\begin{formula}[title={Ecuación ordinaria}]
\[
(x-h)^2+(y-k)^2=r^2.
\]
Centro: $C(h,k)$. Radio: $r$.
\end{formula}

\begin{example}{Centro y radio}{circulo-elementos}
Identifica los elementos de $(x+3)^2+(y-2)^2=25$.
\begin{solution}
\textbf{Paso 1.} Comparamos con $(x-h)^2+(y-k)^2=r^2$.

\textbf{Paso 2.} Como $x+3=x-(-3)$, se tiene $h=-3$; además, $k=2$.

\textbf{Paso 3.} $r^2=25$, por lo que $r=5$. El centro es $(-3,2)$.
\end{solution}
\end{example}

\section{Obtención de la ecuación}

\begin{example}{Centro y un punto}{circulo-punto}
Obtén la circunferencia con centro $C(2,-1)$ que pasa por $P(5,3)$.
\begin{solution}
\textbf{Paso 1.} Calculamos el radio:
\[
r=CP=\sqrt{(5-2)^2+(3+1)^2}=5.
\]
\textbf{Paso 2.} Sustituimos $h=2$, $k=-1$ y $r=5$:
\[
(x-2)^2+(y+1)^2=25.
\]
\end{solution}
\end{example}

\begin{example}{Extremos de un diámetro}{circulo-diametro}
Los extremos de un diámetro son $A(-4,2)$ y $B(6,8)$. Obtén la ecuación.
\begin{solution}
\textbf{Paso 1.} El centro es el punto medio:
\[
C=\left(\frac{-4+6}{2},\frac{2+8}{2}\right)=(1,5).
\]
\textbf{Paso 2.} El radio es la mitad de $AB$:
\[
r=\frac12\sqrt{10^2+6^2}=\sqrt{34}.
\]
\textbf{Paso 3.} La ecuación es
\[
(x-1)^2+(y-5)^2=34.
\]
\end{solution}
\end{example}

\section{Ecuación general}

Al desarrollar la ecuación ordinaria se obtiene una ecuación sin término
$xy$ y con los mismos coeficientes para $x^2$ y $y^2$.

\begin{formula}[title={Ecuación general de la circunferencia}]
\[
x^2+y^2+Dx+Ey+F=0.
\]
Su centro y radio son
\[
C\left(-\frac D2,-\frac E2\right),\qquad
r=\sqrt{\frac{D^2+E^2}{4}-F}.
\]
\end{formula}

Para pasar a forma ordinaria se agrupan los términos en $x$ y en $y$ y se
completan cuadrados usando
\[
u^2+bu=\left(u+\frac b2\right)^2-\left(\frac b2\right)^2.
\]

\begin{example}{Completar cuadrados}{completar-circulo}
Convierte $x^2+y^2-6x+4y-12=0$ a forma ordinaria.
\begin{solution}
\textbf{Paso 1.} Agrupamos y trasladamos la constante:
\[
(x^2-6x)+(y^2+4y)=12.
\]
\textbf{Paso 2.} Añadimos $9$ y $4$ a ambos lados:
\[
(x^2-6x+9)+(y^2+4y+4)=25.
\]
\textbf{Paso 3.} Factorizamos:
\[
(x-3)^2+(y+2)^2=25.
\]
Centro $(3,-2)$ y radio $5$.
\end{solution}
\end{example}

\begin{remark}
Si $r^2>0$ hay una circunferencia real; si $r^2=0$ la gráfica se reduce a un
punto; si $r^2<0$ no existen puntos reales.
\end{remark}

\section{Posiciones relativas}

Para comparar un punto $P$ con una circunferencia se calcula $d=CP$:
\[
d<r:\text{ interior},\qquad d=r:\text{ sobre ella},\qquad d>r:\text{ exterior}.
\]

Para una recta se calcula la distancia $d(C,\ell)$ del centro a la recta:
\[
d<r:\text{ secante},\qquad d=r:\text{ tangente},\qquad d>r:\text{ exterior}.
\]

\begin{example}{Recta y circunferencia}{recta-circulo}
Clasifica la recta $3x+4y-25=0$ respecto de
$(x-1)^2+(y-2)^2=9$.
\begin{solution}
\textbf{Paso 1.} $C(1,2)$ y $r=3$.

\textbf{Paso 2.} Calculamos la distancia a la recta:
\[
d=\frac{|3(1)+4(2)-25|}{\sqrt{3^2+4^2}}=\frac{14}{5}=2.8.
\]
\textbf{Paso 3.} Como $2.8<3$, la recta es secante.
\end{solution}
\end{example}

\section{Circunferencia por tres puntos}

Tres puntos no colineales determinan una circunferencia. Se sustituyen en
$x^2+y^2+Dx+Ey+F=0$ y se resuelve el sistema para $D,E,F$.

\begin{example}{Circunferencia por tres puntos}{tres-puntos-circulo}
Obtén la circunferencia que pasa por $(0,0)$, $(4,0)$ y $(0,6)$.
\begin{solution}
\textbf{Paso 1.} Con $(0,0)$ se obtiene $F=0$.

\textbf{Paso 2.} Con $(4,0)$: $16+4D=0$, así que $D=-4$.

\textbf{Paso 3.} Con $(0,6)$: $36+6E=0$, así que $E=-6$.

\textbf{Paso 4.} La ecuación general es
\[
x^2+y^2-4x-6y=0,
\]
y la ordinaria es $(x-2)^2+(y-3)^2=13$.
\end{solution}
\end{example}

\section{Ejercicios y problemas}

\begin{exercise}{Circunferencia}{ej-circulo}
\begin{enumerate}
  \item Identifica centro y radio de $(x-4)^2+(y+1)^2=36$.
  \item Obtén la circunferencia con centro $(-2,5)$ y radio $4$.
  \item Obtén la ecuación si el centro es $(3,1)$ y pasa por $(-1,4)$.
  \item Convierte $x^2+y^2+8x-10y+16=0$ a forma ordinaria.
  \item Decide si $P(5,2)$ está dentro, sobre o fuera de $x^2+y^2=25$.
  \item Clasifica $x+y-7=0$ respecto de $(x-2)^2+(y-1)^2=8$.
  \item Obtén la circunferencia que pasa por $(1,0)$, $(5,0)$ y $(3,4)$.
\end{enumerate}
\end{exercise}

\begin{problem}{Aplicaciones}{prob-circulo}
\begin{enumerate}
  \item Una zona de cobertura tiene centro $(2,-3)$ y radio $8$ km. Escribe
  su ecuación y decide si el punto $(8,1)$ tiene cobertura.
  \item Determina la ecuación de la circunferencia cuyo diámetro está sobre
  la recta $y=2$ y tiene extremos en $x=-5$ y $x=7$.
  \item Encuentra los valores de $k$ para que $(k,4)$ pertenezca a
  $x^2+y^2-6x-4y-12=0$.
\end{enumerate}
\end{problem}

\section*{Resumen}\addcontentsline{toc}{section}{Resumen}
\begin{properties}
\[
(x-h)^2+(y-k)^2=r^2,
\qquad
x^2+y^2+Dx+Ey+F=0.
\]
\[
C\left(-\frac D2,-\frac E2\right),
\qquad
r=\sqrt{\frac{D^2+E^2}{4}-F}.
\]
\end{properties}
